\documentclass{article}
\usepackage{amsmath}
\begin{document}
\section*{Another duty's past effect can supply present resource context}
Extend the gate proof scaffold to two separately named duties.
The external contract declares one exclusive nonrenewable gate
effect resource: two independently authorized commands do not
automatically merge into one longer closure. Their actual effects
must not overlap. Transfer changes the present operator, not the
shared physical resource or an already committed effect interval.
Both modeled operators may read already settled results within
this fixed gate scope.

For a new duty the context is the current authority joined with
the gate's state of actual occupied intervals, regardless of who
installed them. The content is its independently qualified body
fact and the exact historical settlement of its own semantic
key if present. The mixed decision order is historical read,
live actor check, qualified physical fact, and then compatibility
with other duties' committed intervals. A conflict returns busy
without consuming the requested key; it does not become a renewed
historical permission or an immutable rejection of all later
attempts at that work. This extra resource premise did not appear
in the one-key kernel.

Two native SQLite subprocess histories have the same successor,
qualified clearance and absent own result for duty V. In the first
the original Q worker died before commit; in the second Q has
already installed an exact effect. Nothing in V's own unresolved
content or the new owner's local record distinguishes them.
Forgetting the old owner's live physical lease fails the downstream
compatibility test. The qualified context includes it, rejecting V
both during the initial interval and at its inclusive final tick,
and allowing actual V on the following tick. Old Q still returns
its unchanged recorded result, not extended occupancy. Both modeled
gate identities and the two actual transaction crash cuts execute.
The owner-filtered control installs two conflicting pairs; the
resource-scoped projection installs none and preserves exact Q
history. This comparison is about the declared exclusion
obligation, not a claim of unavoidable physical collision if an
application has different merging authority.

This demonstrates why the implementation cannot simply compose
local handler success bits: an effect surviving on the content
side must be available at the appropriate shared context boundary
for the next duty. Atomicity, one serialized store, stable geometry,
fixed scopes and the actual compatibility predicate are concrete
proof obligations, not consequences of the notation. No general
pasting theorem, distributed agreement, radio or independent
mission acceptance is claimed here.
\end{document}
